\section{Introduction and Companies Selected}

As the subject of this study, I chose three companies more or less related to the
transportation industry, which for various reasons have had difficulties in recent
years, especially environment-related, to publish good business figures. The two
aircraft manufacturers \textbf{Boeing} and \textbf{Airbus}, which share the market
completely between themselves, were particularly affected by the lockdown during the
last pandemic and the shutdown. However, vehicle manufacturers such as market leader
\textbf{Toyota} are also facing changing market demand. In view of these challenges, it
is interesting to conduct an analysis of future performance and to compare buy-and-hold
strategies to our trained agent.

The following descriptive statistics are used to evaluate the performance of the shares:
\begin{itemize}
  \item share price over the last 10 years,
  \item price-earnings ratio (P/E ratio),
  \item price-to-sales ratio (P/S ratio),
  \item log-returns.
\end{itemize}

\paragraph{Price-earnings ratio (P/E ratio).}
The price/earnings ratio is the number of years in which the company would have earned
its stock market value if profits had remained constant. In the calculation, the current
share or stock market price is divided by the company's earnings per share. The
company's earnings per share are calculated by dividing the company's earnings by the
total number of shares. The common and relevant sites give a guideline value between 12
and 15 for the interpretation of the P/E ratio. If the results are below this value, a
share tends to be considered cheap; if the results are above 15, the share may tend to be
considered more expensive. The formula is
\begin{equation}
  P/E = \frac{\text{share price}}{\text{earnings per share}}.
\end{equation}

\paragraph{Price-to-sales ratio (P/S ratio).}
The price-to-sales ratio (P/S ratio) is a share ratio that shows the market value of a
share in relation to its annual sales. The lower the P/S compared to the values of other
companies in the same industry, the more favorable the respective share is. The common
and relevant sites give a guideline value above 1 and below 1 for the interpretation of
the P/S ratio. If the results are below this value, a share tends to be considered cheap;
if the results are above 1, the share may tend to be considered more expensive. It can be
calculated as follows:
\begin{equation}
  P/S_{total} = \frac{\text{sales}}{\text{market capitalisation}}.
\end{equation}

\paragraph{Log-returns.}
The log return can be calculated by simply taking a natural logarithm of the current
price $(p_1)$ divided by the initial price $(p_0)$:
\begin{equation}
  r_1 = \ln\!\left(\frac{p_1}{p_0}\right).
\end{equation}
